	\begin{UseCase}{CU6}{Gestionar Términos de Glosario}{
	Un término del glosario es una palabra que junto con su descripción ayudará a comprender  aspectos del negocio del sistema.\\
	
	Este caso de uso permite al colaborador (ya sea \hyperlink{jefe}{Líder} o \hyperlink{analista}{Analista}) visualizar en una tabla el listado de  todos los términos del glosario que componen el proyecto en el que se está operando y solicitar el registro de uno nuevo. En cada módulo existente se podrán realizar diferentes acciones como: editar la descripción de algún término previamente registrado (para actualizar nombre o descripción) o eliminar un término (en caso de que este ya no tenga una razón de ser dentro del sistema).\\
	}
	\UCitem{Actor}{\hyperlink{jefe}{Líder de análisis}, \hyperlink{analista}{Analista}}
	\UCitem{Propósito}{Proporcionar al actor un mecanismo para llevar el control de los términos de un glosario de un proyecto.}
	\UCitem{Entradas}{Ninguna}
	\UCitem{Salidas}{\begin{itemize}
			\item \cdtRef{proyectoEntidad:claveProyecto}{Clave del proyecto}: Lo obtiene el sistema.
			\item \cdtRef{proyectoEntidad:nombreProyecto}{Nombre del proyecto}: Lo obtiene el sistema.
			\item \cdtRef{terminoGLSEntidad}{Término del glosario}: Tabla que muestra \cdtRef{terminoGLSEntidad:claveTGLS}{Clave}, y el \cdtRef{moduloEntidad:nombreTGLS}{Nombre} de todos los términos registrados de un proyecto.
			\item \cdtIdRef{MSG2}{No existe información}: Se muestra en la pantalla \IUref{IU11A}{Gestionar Términos de Glosario} cuando no existen términos registrados.
	\end{itemize}}
	
	\UCitem{Precondiciones}{Que exista al menos un proyecto registrado.}
	\UCitem{Postcondiciones}{Ninguna}
	\UCitem{Errores}{Ninguno}
	\UCitem{Tipo}{Primario}
\end{UseCase}
%--------------------------------------
\begin{UCtrayectoria}
	\UCpaso[\UCactor] Solicita gestionar los términos seleccionando la opción ''Glosario'' del menú \IUref{MN3}{Menú de Proyecto}.
	\UCpaso[\UCsist] Obtiene la información de los términos registrados del proyecto seleccionado. \hyperlink{CU6:TAA}{[Trayectoria A]}
	\UCpaso[\UCsist] Ordena los términos alfabéticamente basándose en nombre de los mismos.
	\UCpaso[\UCsist] Muestra la información de los términos en la pantalla \IUref{IU11}{Gestionar Términos de glosario} y las operaciones disponibles de acuerdo a la regla de negocio \BRref{RN15}{Operaciones disponibles}.\label{CU6-P4}
	\UCpaso[\UCactor] Gestiona los términos a través de los botones: \IUbutton{Registrar}, \editar, \eliminar y \raisebox{-1mm}{\includegraphics[height=11pt]{images/Iconos/consultar}}. 
\end{UCtrayectoria}		
%--------------------------------------
\hypertarget{CU6:TAA}{\textbf{Trayectoria alternativa A}}\\
\noindent \textbf{Condición:} No existen registros de términos.
\begin{enumerate}
	\UCpaso[\UCsist] Muestra el mensaje \cdtIdRef{MSG2}{No existe información} en la pantalla \IUref{IU11A}{Gestionar Términos de Glosario} para indicar que no hay registros de términos para mostrar. \label{CU6-TA1}
	\UCpaso[\UCactor] Gestiona los términos a través del botón: \IUbutton{Registrar}. 
	\item[- -] - - {\em {Fin del caso de uso}}.%
\end{enumerate}

%--------------------------------------

\subsubsection{Puntos de extensión}

\UCExtenssionPoint{El actor requiere registrar un término.}{Presionando el botón \IUbutton{Registrar} del paso \ref{CU6-P4} de la trayectoria principal o del paso \ref{CU6-TA1} de la trayectoria alternativa A.}{\UCref{CU6.1}{Registrar Término}}
\UCExtenssionPoint{El actor requiere modificar un término.}{Presionando el icono \editar del paso \ref{CU6-P4} de la trayectoria principal.}{\UCref{CU6.2}{Modificar Término}}
\UCExtenssionPoint{El actor requiere eliminar un término.}{Presionando el icono \eliminar del paso \ref{CU6-P4} de la trayectoria principal.}{\UCref{CU6.3}{Eliminar Término}}
\UCExtenssionPoint{El actor requiere consultar un término.}{Presionando el icono \raisebox{-1mm}{\includegraphics[height=11pt]{images/Iconos/consultar}} del paso \ref{CU6-P4} de la trayectoria principal.}{\UCref{CU6.4}{Consultar Término}}